\documentclass[<<options>>]{apmcmthesis}
\tihao{<<A>>}
\baominghao{apmcm2699001}   % <-- 队号，提交前必须填；必须 apmcm+数字，否则页眉正则配不上
\begin{document}
% 官方格式（2017-2023 每年 Essay format p.67 起；2026 扫描件 p3/p4 同义）：
%   第 1 页 = Summary Sheet（摘要页），第 2 页 = Contents，第 3 页起正文。
% 英文赛项**没有** Control Sheet（承诺书是中文赛项那一套）。
% 实测坑（2018A 练习稿）：HEAD 原来先 \makecontrolsheet 再 \maketitle，承诺书占第 1 页，
% review_check 的 C4「第一页必须是摘要页」直接判红。而且 \makecontrolsheet 会把 \maketitle
% \let 成 relax，所以删一行不够 —— 两行都要去掉，直接从 \begin{abstract} 起。
\thispagestyle{empty}
\begin{abstract}
\textbf{Background and Task} <<one sentence with a number>>. This paper addresses
<<N>> problems: <<task 1>>; <<task 2>>; <<task 3>>. <<N>> models are developed.

\textbf{Modeling Framework} <<Q1 model>>, based on which <<Q2 model>>, and extending
to <<Q3 model>>.

\textbf{Results per Question} <<Q1 number + conclusion>>; <<Q2 number + conclusion>>;
<<Q3 number + conclusion>>.

\textbf{Validation and Limitation} <<validation 1 with number>>; <<validation 2 with
number>>. The main limitation is that <<one concrete limitation with a number>>.

\keywords{<<model names in the order they appear above>>}
\end{abstract}

\newpage
\tableofcontents
\newpage
% 必须是 mainmatterstyle：它才是 \fancyhead 里带 "Team # ... Page n of N" 的那个。
% 用 plain 会丢掉页眉，gauntlet 的 R-V02 直接判 body=0。
\pagestyle{mainmatterstyle}
\setcounter{page}{1}

\section{Introduction}
\subsection{Background}
<<why this matters now, one paragraph>>
\subsection{Restatement of the Problems}
<<restate each question in your own variables>>
\subsection{Our Work}
<<what this paper delivers, one line per question>>
\section{Problem Analysis}
<<the mechanism, the key tension, what you simplify and why>>
\section{Assumptions and Justifications}
\begin{enumerate}
  \item <<assumption>> --- <<why>>; used by Section~<<x>>.
  \item <<assumption>> --- <<why>>; used by Section~<<x>>.
\end{enumerate}
\section{Notations}
\begin{table}[htbp]\centering\caption{Notations}\label{tab:notation}
\begin{tabular}{@{}lll@{}}\toprule Symbol & Meaning & Unit \\ \midrule
$<<sym>>$ & <<meaning>> & <<unit>> \\ \bottomrule\end{tabular}\end{table}
\section{Question 1: <<What it does>>}
\subsection{Image Pre-processing}
<<fill>>
\subsection{<<Discriminant>> Modeling}
<<fill>>
\subsection{Solution and Results}
<<fill>>
\subsection{Test Analysis}
<<fill>>
\subsection{Short Summary}
<<fill>>
\section{Question 2: <<What it does>>}
\subsection{Pre-processing}
<<fill>>
\subsection{Model}
<<fill>>
\subsection{Solution and Results}
<<fill>>
\subsection{Test Analysis}
<<fill>>
\subsection{Short Summary}
<<fill>>
\section{Sensitivity Analysis and Error Analysis}
<<perturbation / resampling / variance decomposition; report the unfavourable half>>
\section{Conclusion}
<<per-question summary, each pointing at a number>>
\section{Strengths and Weaknesses}
\subsection{Strengths}
\begin{itemize}
  \item <<strength 1, points to a result>>
  \item <<strength 2>>
  \item <<strength 3>>
\end{itemize}
\subsection{Weaknesses}
\begin{itemize}
  \item <<weakness 1, with a number>>
  \item <<weakness 2, with a number>>
  \item <<weakness 3, with a number>>
\end{itemize}

\begin{thebibliography}{99}
\bibitem{ref1} <<author>>. <<title>>[J]. <<journal>>, <<year>>, <<vol>>(<<no>>): <<pages>>.
\bibitem{ref2} <<author>>. <<book>>[M]. <<place>>: <<publisher>>, <<year>>.
\bibitem{ref3} <<institution>>. <<report>>[R]. <<year>>.
\bibitem{ref4} <<institution>>. <<dataset>>[DB/OL]. <<url>>, accessed <<YYYY-MM-DD>>.
\end{thebibliography}

\appendix
% 必须用 \section{Appendix A: Code}：apmcmthesis.cls 里 \section 被改成罗马数字（I. ），
% \appendix 不会把它变成 "Appendix A"。而 gauntlet 认正文首页的锚点正是字面 "Appendix A"。
\section{Appendix A: Code}
% 先把 code/main.py 写出来，再把下一行取消注释。模板默认不启用：
% 缺文件时 listings 抛 "! Package Listings Error: File `code/main(.py)' not found"
% + "! Emergency stop."，exit=1，且 lastpage 跳号让页眉变成 "Page N of ??"。
%\lstinputlisting{code/main.py}
<<paste the full source of every script that produced a number, one appendix per script>>
\end{document}
